\documentclass[UTF8,fontset=none,11pt]{ctexart}
\usepackage[a4paper,margin=2.05cm]{geometry}
\usepackage{amsmath,amssymb}
\usepackage{enumitem}
\usepackage[table]{xcolor}
\usepackage{booktabs}
\usepackage{array}
\usepackage{longtable}
\usepackage{tikz}
\usepackage{graphicx}
\usepackage{hyperref}
\usetikzlibrary{positioning,arrows.meta,shapes.geometric,fit,backgrounds,calc}

\setCJKmainfont[AutoFakeSlant=0.18]{PingFang SC}
\setCJKsansfont[AutoFakeSlant=0.18]{PingFang SC}
\setCJKmonofont[AutoFakeSlant=0.18]{PingFang SC}

\hypersetup{
  colorlinks=true,
  linkcolor=blue!55!black,
  urlcolor=blue!55!black,
  citecolor=blue!55!black,
  pdftitle={Syntax Analysis Bottom-Up 中文学习笔记},
  pdfauthor={Trae AI}
}

\setlength{\parindent}{0pt}
\setlength{\parskip}{0.55em}
\setlist[itemize]{leftmargin=1.8em,itemsep=0.22em,topsep=0.25em}
\setlist[enumerate]{leftmargin=2.0em,itemsep=0.28em,topsep=0.3em}
\renewcommand{\arraystretch}{1.18}
\setlength{\emergencystretch}{3em}
\sloppy

\definecolor{defrule}{RGB}{25,92,140}
\definecolor{noterule}{RGB}{120,80,15}
\definecolor{warnrule}{RGB}{160,50,45}
\definecolor{softblue}{RGB}{235,246,255}
\definecolor{softyellow}{RGB}{255,249,230}
\definecolor{softred}{RGB}{255,239,239}
\definecolor{graytext}{RGB}{95,95,95}
\definecolor{GRAYTEXT}{RGB}{95,95,95}
\newcolumntype{L}[1]{>{\raggedright\arraybackslash}p{#1}}

\newenvironment{defbox}[1][]{%
  \par\medskip\noindent{\color{defrule}\rule{\linewidth}{1.1pt}}\par\nopagebreak
  \noindent{\bfseries\color{defrule}#1}\par\nopagebreak\smallskip}%
  {\par\nopagebreak\smallskip\noindent{\color{defrule}\rule{\linewidth}{0.4pt}}\par\medskip}
\newenvironment{notebox}[1][]{%
  \par\medskip\noindent{\color{noterule}\rule{\linewidth}{1.1pt}}\par\nopagebreak
  \noindent{\bfseries\color{noterule}#1}\par\nopagebreak\smallskip}%
  {\par\nopagebreak\smallskip\noindent{\color{noterule}\rule{\linewidth}{0.4pt}}\par\medskip}
\newenvironment{warnbox}[1][]{%
  \par\medskip\noindent{\color{warnrule}\rule{\linewidth}{1.1pt}}\par\nopagebreak
  \noindent{\bfseries\color{warnrule}#1}\par\nopagebreak\smallskip}%
  {\par\nopagebreak\smallskip\noindent{\color{warnrule}\rule{\linewidth}{0.4pt}}\par\medskip}

\newcommand{\pg}[1]{\texorpdfstring{{\small\color{graytext}\,[p#1]}}{ [p#1]}}
\newcommand{\eps}{\varepsilon}
\newcommand{\FIRST}{\ensuremath{\mathrm{FIRST}}}
\newcommand{\FOLLOW}{\ensuremath{\mathrm{FOLLOW}}}
\newcommand{\ACTION}{\ensuremath{\mathrm{ACTION}}}
\newcommand{\GOTO}{\ensuremath{\mathrm{GOTO}}}
\newcommand{\CLOSURE}{\ensuremath{\mathrm{CLOSURE}}}
\newcommand{\LR}{\ensuremath{\mathrm{LR}}}
\newcommand{\SLR}{\ensuremath{\mathrm{SLR}}}
\newcommand{\LALR}{\ensuremath{\mathrm{LALR}}}

\title{\bfseries Lecture 5: Syntax Analysis Bottom-Up\\中文学习笔记}
\author{基于课件 \texttt{5.Syntax\_Analysis\_Bottom\_Up.pdf} 整理}
\date{\today}

\begin{document}
\maketitle
\tableofcontents
\newpage

\section*{阅读导引}
\addcontentsline{toc}{section}{阅读导引}

本讲围绕语法分析的自底向上路线展开，核心是 LR 系列分析器。前半部分讲如何从移入-规约分析走到 LR(0) 自动机和解析表；中间依次引入 SLR(1)、LR(1)、LALR(1) 来解决 LR(0) 过弱和 LR(1) 表过大的问题；后半部分复习前几讲语法分析内容，并给出练习题讲解。

\begin{notebox}[一条主线]
自底向上分析的根本问题是：栈顶什么时候已经形成句柄，应该用哪条产生式规约。LR 系列的解决思路是把“当前已经识别到哪里”编码成状态，再用 \ACTION/\GOTO 表驱动移入和规约。
\end{notebox}

\section{课程总览：Bottom-up 到 LR Family \pg{1--3}}

课件开头给出全局路线图：

\begin{center}
\resizebox{0.98\linewidth}{!}{%
\begin{tikzpicture}[font=\small, node distance=0.9cm, >=Latex]
  \tikzstyle{n}=[draw,rounded corners,align=center,minimum height=0.85cm,fill=softblue]
  \tikzstyle{hi}=[draw,rounded corners,align=center,minimum height=0.85cm,fill=softyellow]
  \node[n] (bu) {Bottom-up\\parser};
  \node[n,right=1.0cm of bu] (sr) {Shift-Reduce\\句柄/活前缀/冲突};
  \node[n,right=1.0cm of sr] (lr0) {LR(0)\\items + DFA};
  \node[hi,above right=0.55cm and 0.9cm of lr0] (slr) {SLR(1)\\FOLLOW 限制规约};
  \node[hi,right=1.0cm of lr0] (lr1) {LR(1)\\item + lookahead};
  \node[hi,below right=0.55cm and 0.9cm of lr0] (lalr) {LALR(1)\\合并同 core 状态};
  \node[n,right=1.1cm of lr1] (table) {ACTION/GOTO\\parser table};
  \draw[->] (bu) -- (sr);
  \draw[->] (sr) -- (lr0);
  \draw[->] (lr0) -- (slr);
  \draw[->] (lr0) -- (lr1);
  \draw[->] (lr1) -- (lalr);
  \draw[->] (slr) -- (table);
  \draw[->] (lr1) -- (table);
  \draw[->] (lalr) -- (table);
\end{tikzpicture}
}
\end{center}

LR 系列的几个变体可以先这样理解：

\begin{longtable}{L{0.16\linewidth}L{0.26\linewidth}L{0.28\linewidth}L{0.20\linewidth}}
\toprule
类型 & 使用的信息 & 规约策略 & 特点 \\
\midrule
LR(0) & 只有 LR(0) item 和栈状态 & 完成项出现就规约 & 最简单、最弱，容易冲突 \\
SLR(1) & LR(0) 状态 + FOLLOW 集 & 下一个输入符号在 \(\FOLLOW(A)\) 中才用 \(A\to\alpha\) 规约 & 状态数与 LR(0) 相同，能力更强 \\
LR(1) & LR(0) item + 精确 lookahead & 下一个输入符号等于 item 的 lookahead 才规约 & 能力强，但状态数和表可能很大 \\
LALR(1) & LR(1) lookahead + 合并同 core 状态 & 保留 lookahead，但减少状态数 & 实用折中，YACC/Bison 常用路线 \\
\bottomrule
\end{longtable}

\section{自底向上分析与移入-规约 \pg{4--9}}

\subsection{Bottom-up Parsing 的基本思想 \pg{4--5}}

\begin{defbox}[定义：自底向上分析]
自底向上分析（bottom-up parsing）从输入串出发，逐步把某些子串替换为产生式左部的非终结符，直到整个输入被规约为开始符号。它等价于“最右推导的逆过程”，语法树是从叶子结点向根结点构造出来的。
\end{defbox}

通俗地说，Top-down 是“从开始符号猜出输入串”，Bottom-up 是“从输入串往回合成开始符号”。例如表达式文法：
\[
E\to E+T\mid T,\qquad T\to T*F\mid F,\qquad F\to (E)\mid id
\]
对输入 \(id*id\)，一种最右推导为
\[
E \Rightarrow T \Rightarrow T*F \Rightarrow T*id \Rightarrow F*id \Rightarrow id*id.
\]
自底向上分析反过来做：
\[
id*id \Rightarrow F*id \Rightarrow T*id \Rightarrow T*F \Rightarrow T \Rightarrow E.
\]

自底向上分析相对自顶向下分析更强：

\begin{itemize}
  \item 不要求对文法做左因子提取（left factoring）。
  \item 可以处理左递归（left recursion），而递归下降分析直接遇到左递归会无限递归。
  \item 代价是手写分析器很困难，通常由工具根据文法自动生成。
\end{itemize}

\subsection{Shift-Reduce Parser 的四种动作 \pg{6--7}}

\begin{defbox}[定义：移入-规约分析]
Shift-reduce parsing 是自底向上分析的一种形式：栈保存已经读入的文法符号，输入缓冲区保存尚未扫描的输入串。分析器从左到右扫描输入，反复执行移入和规约。
\end{defbox}

\begin{longtable}{L{0.17\linewidth}L{0.73\linewidth}}
\toprule
动作 & 含义 \\
\midrule
Shift 移入 & 把下一个输入符号压入栈顶。通常表示句柄还没有完整出现在栈顶，需要继续读输入。\\
Reduce 规约 & 栈顶某段符号串匹配某条产生式右部 \(\beta\)，就把它替换为产生式左部 \(A\)。也就是执行 \(\beta \Rightarrow A\) 的逆向合成。\\
Accept 接受 & 输入已经成功规约为开始符号，分析完成。\\
Error 报错 & 当前状态和输入符号无合法动作，说明语法错误，需要报错或恢复。\\
\bottomrule
\end{longtable}

\subsection{例子：分析 \texorpdfstring{\(abbcde\)}{abbcde} \pg{8--9}}

给定文法：
\[
S\to aAcBe,\qquad A\to b,\qquad A\to Ab,\qquad B\to d.
\]
输入为 \(abbcde\)。课件展示的分析过程如下：

\begin{center}
\small
\begin{tabular}{lll}
\toprule
栈 & 输入 & 动作 \\
\midrule
\(\$\) & \(abbcde\$\) & Shift \\
\(\$a\) & \(bbcde\$\) & Shift \\
\(\$ab\) & \(bcde\$\) & Reduce \(A\to b\) \\
\(\$aA\) & \(bcde\$\) & Shift \\
\(\$aAb\) & \(cde\$\) & Reduce \(A\to Ab\) \\
\(\$aA\) & \(cde\$\) & Shift \\
\(\$aAc\) & \(de\$\) & Shift \\
\(\$aAcd\) & \(e\$\) & Reduce \(B\to d\) \\
\(\$aAcB\) & \(e\$\) & Shift \\
\(\$aAcBe\) & \(\$\) & Reduce \(S\to aAcBe\) \\
\(\$S\) & \(\$\) & Accept \\
\bottomrule
\end{tabular}
\end{center}

这个例子引出自底向上分析的两个核心决策：

\begin{itemize}
  \item 什么时候应该继续 Shift，什么时候应该 Reduce？
  \item Reduce 时到底应该用哪条产生式？
\end{itemize}

例如栈顶出现 \(b\) 时，既可能用 \(A\to b\) 规约，也可能继续读入，等待形成 \(Ab\) 再用 \(A\to Ab\) 规约。解决这个问题需要引入“句柄”。

\section{句柄、短语、直接短语与活前缀 \pg{10--27}}

\subsection{右句型、短语、直接短语、句柄 \pg{10--12}}

\begin{defbox}[右句型]
右句型（right-sentential form）是在最右推导过程中出现的句型。自底向上分析做的是最右推导的逆过程，因此它关心当前输入对应的右句型。
\end{defbox}

\begin{defbox}[短语、直接短语、句柄]
设 \(\alpha_1A\alpha_2\) 是文法 \(G(S)\) 的一个右句型，且
\[
S\Rightarrow^* \alpha_1A\alpha_2 \Rightarrow \alpha_1\beta\alpha_2.
\]
则 \(\beta\) 是该句型的一个短语（phrase）。若这一步 \(A\Rightarrow \beta\) 对应语法树中只有父子两代的子树，则 \(\beta\) 是直接短语（simple phrase）。句柄（handle）是右句型中最左的直接短语。
\end{defbox}

通俗解释：

\begin{itemize}
  \item 短语：语法树中某个子树的所有叶子拼成的串。
  \item 直接短语：该子树只有“父结点 - 子结点”两层，说明它刚好对应一条产生式右部。
  \item 句柄：当前最应该先规约的那一个直接短语。自底向上分析每一步都应该规约句柄。
\end{itemize}

对文法
\[
E\to T\mid E+T,\qquad T\to F\mid T*F,\qquad F\to (E)\mid a\mid b\mid c
\]
和句型 \(a*b+c\)，课件给出的结果为：

\begin{center}
\small
\begin{tabular}{ll}
\toprule
类别 & 符号串 \\
\midrule
短语 & \(a*b+c,\ a*b,\ a,\ b,\ c\) \\
直接短语 & \(a,\ b,\ c\) \\
句柄 & \(a\) \\
\bottomrule
\end{tabular}
\end{center}

注意 \(b+c\)、\(*b\)、\(+\) 不是短语，因为它们不能对应语法树中某个完整子树的叶子序列。

\subsection{句柄规约示例 \pg{13--24}}

课件用句型 \(T+T*F+id\) 展示句柄如何逐步被规约。对应表达式文法：
\[
E\to E+T\mid T,\qquad T\to T*F\mid F,\qquad F\to (E)\mid id.
\]
该句型中的短语包括：
\[
T+T*F+id,\quad T+T*F,\quad T,\quad T*F,\quad id.
\]
直接短语包括左侧的 \(T\)、\(T*F\) 中的可规约部分，以及 \(id\) 等；每一步真正规约的是最左直接短语，即句柄。

核心规约序列可以概括为：
\[
T+T*F+id
\Rightarrow E+T*F+id
\Rightarrow E+T+id
\Rightarrow E+id
\Rightarrow E+F
\Rightarrow E+T
\Rightarrow E.
\]

在 \(abbcde\) 的例子中，栈为 \(\$aAb\)、输入为 \(cde\$\) 时，句柄是 \(Ab\)，因此应选择 \(A\to Ab\) 规约，而不是错误地只把 \(b\) 按 \(A\to b\) 规约。

\subsection{为什么句柄总在栈顶 \pg{25}}

句柄可能一开始还在输入缓冲区中，没有完全进入栈；但分析器会继续 Shift，直到句柄完整出现在栈顶。一旦句柄在栈顶，分析器会立即 Reduce。因此句柄不会“长期卡在栈中间”。

这给出通用移入-规约策略：

\begin{itemize}
  \item 若栈顶没有句柄，Shift。
  \item 若栈顶已经形成句柄，Reduce 成相应非终结符。
  \item 可通过表自动化实现上述决策。
\end{itemize}

\subsection{活前缀 \pg{26--27}}

\begin{defbox}[活前缀]
活前缀（viable prefix）是某个右句型的前缀，并且该前缀没有越过该右句型最右句柄的右端。
\end{defbox}

直观上，活前缀是“仍然有希望成为正确规约路径的栈内容”。LR 分析过程中，栈中符号串始终是某个右句型的活前缀，因此分析器不会偏离正确方向。

课件例子：
\[
S\Rightarrow aEb\Rightarrow aaEb.
\]
若句柄为 \(aE\)，则活前缀包括 \(a, aa, aaE\)，但不包括 \(aaEb\)，因为 \(aaEb\) 已经越过句柄右端。

另一个例子中，右句型 \(aAcde\) 可被分在栈和输入之间；栈内容 \(aA\)、\(aAc\)、\(aAcd\) 都是活前缀。

\section{冲突与自底向上分析性质 \pg{28--31}}

\subsection{冲突来自哪里 \pg{28--30}}

考虑二义表达式文法：
\[
E\to E*E\mid E+E\mid (E)\mid id.
\]
对输入 \(id*id+id\)，当栈中已经有 \(E*E\)、下一个输入符号是 \(+\) 时，有两种选择：

\begin{itemize}
  \item 先 Reduce \(E\to E*E\)，得到 \((id*id)+id\)，体现 \(*\) 优先于 \(+\)。
  \item 先 Shift \(+\)，之后再规约，得到 \(id*(id+id)\)。
\end{itemize}

文法没有说明优先级和结合性，所以两个选择都看似合理。这就是冲突。

\begin{defbox}[两类冲突]
Shift-reduce conflict：同一状态下既可以移入，也可以规约。Reduce-reduce conflict：同一状态下可用两条或多条产生式规约。
\end{defbox}

大多数冲突来自二义文法，但也存在非二义文法无法被某类弱分析器处理的情况。处理冲突的常见方式包括：

\begin{itemize}
  \item 改写文法，把运算符优先级和结合性编码进去。
  \item 对特殊结构手动消除二义性，例如 dangling-else。
  \item 使用更强的分析器类别，例如从 LR(0) 提升到 SLR(1)、LR(1) 或 LALR(1)。
\end{itemize}

\subsection{Bottom-up Parsing 的性质 \pg{31}}

课件总结了三点：

\begin{itemize}
  \item 句柄总会出现在栈顶，这说明用栈实现移入-规约是合理的。
  \item 栈顶没有句柄就移入，栈顶有句柄就规约。
  \item 自动化的关键是构造表，让分析器知道每个状态和输入下应采取什么动作。
\end{itemize}

\section{LR Parser 与 ACTION/GOTO 表 \pg{32--62}}

\subsection{LR(k) Parser 的含义 \pg{33--34}}

课件列出多种自底向上分析器：simple precedence parser、operator precedence parser、recursive ascent parser、LR family parser 等。本讲主要关注 LR 系列。

\begin{defbox}[LR(k)]
LR(k) 是 LR 系列分析器。L 表示从左到右扫描输入（left-to-right scanning），R 表示反向构造最右推导（rightmost derivation in reverse），k 表示做语法决策时向前看的输入符号数量。
\end{defbox}

与 LL(k) 相比：

\begin{itemize}
  \item 二者都可达到输入长度线性时间和线性空间。
  \item 二者都能由工具根据文法自动生成，例如 YACC、Bison。
  \item LR 比 LL 更强，不需要消除左递归，也不需要左因子提取。
  \item LR 冲突调试更复杂，因为冲突发生在自动机状态和表项中，不像 LL 那样直接对应某个递归函数选择。
\end{itemize}

\subsection{LR Parser 的栈与右句型 \pg{35}}

LR parser 的栈保存状态序列：
\[
S_0S_1\cdots S_m,
\]
其中 \(S_m\) 是栈顶状态。每个文法符号 \(X_i\) 会关联一个状态 \(S_i\)。若输入串属于该语言，则“栈中符号 + 未读输入”构成某个右句型：
\[
X_1X_2\cdots X_m a_i a_{i+1}\cdots a_n.
\]

分析器根据栈顶状态 \(S_m\) 和当前输入符号 \(a_i\) 查表：
\[
\ACTION[S_m,a_i].
\]

\subsection{ACTION 表和 GOTO 表 \pg{37--41}}

LR 解析表由两个部分组成：

\begin{longtable}{L{0.18\linewidth}L{0.72\linewidth}}
\toprule
表 & 功能 \\
\midrule
\ACTION 表 & 按“状态 + 终结符”索引，决定移入、规约、接受或报错。\\
\GOTO 表 & 按“状态 + 非终结符”索引，在规约后决定新的栈顶状态。\\
\bottomrule
\end{longtable}

\ACTION 的四种结果如下：

\begin{longtable}{L{0.16\linewidth}L{0.74\linewidth}}
\toprule
表项 & 执行方式 \\
\midrule
\(s_x\) & Shift 当前输入符号，并把状态 \(x\) 压栈。\\
\(r_y\) & Reduce。若第 \(y\) 条产生式为 \(A\to\beta\)，弹出 \(|\beta|\) 个符号及其状态，压入 \(A\)，再查 \(\GOTO\) 压入新状态。\\
acc & Accept，语法分析成功。\\
空表项 & Error，输入不符合文法或需要错误恢复。\\
\bottomrule
\end{longtable}

规约后的关键动作是：
\[
S_j=\GOTO[S_i,A],
\]
其中 \(S_i\) 是弹出句柄后新的栈顶状态，\(A\) 是规约得到的非终结符。

\subsection{表驱动 LR 程序 \pg{42--44}}

LR 程序的伪代码可以概括为：

\begin{enumerate}
  \item 初始时把 \(S_0\) 压栈，输入缓冲区为 \(\omega\$\)。
  \item 令 \(s\) 为栈顶状态，\(a\) 为当前输入符号。
  \item 若 \(\ACTION[s,a]=s_t\)，移入 \(a\) 并压入状态 \(t\)。
  \item 若 \(\ACTION[s,a]=r_k\)，设第 \(k\) 条产生式为 \(A\to\beta\)，弹出 \(|\beta|\) 个符号和状态，压入 \(A\)，再压入 \(\GOTO[s',A]\)，其中 \(s'\) 是弹出后的栈顶状态。
  \item 若 \(\ACTION[s,a]=acc\)，接受。
  \item 否则报错。
\end{enumerate}

\subsection{示例：用 LR 表分析 \texorpdfstring{\(1+1\)}{1+1} \pg{45--58}}

文法为：
\[
(1)\ E\to E*B,\quad
(2)\ E\to E+B,\quad
(3)\ E\to B,\quad
(4)\ B\to 0,\quad
(5)\ B\to 1.
\]
课件给出解析表，并用它分析输入 \(1+1\)。关键步骤如下：

\begin{center}
\small
\begin{tabular}{llll}
\toprule
状态栈 & 符号栈 & 输入 & 动作/跳转 \\
\midrule
0 & \(\$\) & \(1+1\$\) & \(s2\) \\
0 2 & \(\$1\) & \(+1\$\) & \(r5: B\to 1,\ \GOTO(0,B)=4\) \\
0 4 & \(\$B\) & \(+1\$\) & \(r3: E\to B,\ \GOTO(0,E)=3\) \\
0 3 & \(\$E\) & \(+1\$\) & \(s6\) \\
0 3 6 & \(\$E+\) & \(1\$\) & \(s2\) \\
0 3 6 2 & \(\$E+1\) & \(\$\) & \(r5: B\to 1,\ \GOTO(6,B)=8\) \\
0 3 6 8 & \(\$E+B\) & \(\$\) & \(r2: E\to E+B,\ \GOTO(0,E)=3\) \\
0 3 & \(\$E\) & \(\$\) & \(acc\) \\
\bottomrule
\end{tabular}
\end{center}

这个例子说明：LR 分析器本身只负责查表执行动作；真正困难的是构造可靠的解析表，让每个状态和输入符号都能得到正确的 \ACTION/\GOTO 动作。

\subsection{解析表类型概览 \pg{59--62}}

课件将解析表构造分为几类：

\begin{itemize}
  \item LR(0)：只看栈状态，不看任何输入展望符号。简单但太弱。
  \item SLR(1)：在 LR(0) 状态基础上，用 FIRST/FOLLOW 信息给规约加限制。
  \item LALR(1)：常用变体，用更精细的 lookahead 分析，但仍使用接近 LR(0) 的状态规模。
  \item LR(1)：精确的一符号 lookahead，算法和表更复杂。
  \item LR(k)：使用 \(k\) 个展望符号，表更大。
\end{itemize}

\section{LR(0)：项目、自动机与解析表 \pg{63--95}}

\subsection{状态为什么需要 item \pg{63--66}}

问题：当栈顶是 \(T\)、下一个输入符号是 \(*\) 时，分析器如何知道不能立刻把 \(T\) 规约成 \(E\)，而应该先移入 \(*\)？

LR 的答案是维护状态。状态表示“分析已经走到哪里”，由一组 item 构成。

\begin{defbox}[LR(0) item]
LR(0) item 是在产生式右部某个位置加入点号的产生式。对 \(A\to XYZ\)，有四个 item：
\[
A\to \cdot XYZ,\quad A\to X\cdot YZ,\quad A\to XY\cdot Z,\quad A\to XYZ\cdot.
\]
若产生式为 \(A\to\eps\)，只有一个 item：\(A\to\cdot\)。
\end{defbox}

点号的意义：

\begin{itemize}
  \item \(A\to\cdot XYZ\)：希望接下来看到能由 \(XYZ\) 推导出的串。
  \item \(A\to X\cdot YZ\)：已经识别了 \(X\)，接下来期望 \(YZ\)。
  \item \(A\to XYZ\cdot\)：已经识别了完整右部 \(XYZ\)，可能可以规约为 \(A\)。
\end{itemize}

一个状态通常不是单个 item，而是一组 item，称为配置集（configuration set）。例如如果有 item \(A\to X\cdot YZ\)，且 \(Y\to u\mid w\)，则当前状态也应该包含 \(Y\to\cdot u\)、\(Y\to\cdot w\)，因为接下来可能要识别 \(Y\)。

\subsection{增广文法 \pg{67}}

\begin{defbox}[增广文法]
为了让自动机有明确的接受状态，给原文法 \(G\) 增加新的开始符号 \(S'\) 和产生式 \(S'\to S\)，得到增广文法 \(G'\)。当状态中出现 \(S'\to S\cdot\) 且输入为 \(\$\) 时，分析接受。
\end{defbox}

例如原文法：
\[
E\to E*B\mid E+B\mid B,\qquad B\to 0\mid 1
\]
增广后为：
\[
S\to E,\qquad E\to E*B\mid E+B\mid B,\qquad B\to 0\mid 1.
\]

\subsection{CLOSURE 与 GOTO \pg{68--74}}

\begin{defbox}[CLOSURE]
若 \(I\) 是 item 集，\(\CLOSURE(I)\) 按以下规则构造：
\begin{enumerate}
  \item 初始把 \(I\) 中所有 item 加入闭包。
  \item 若闭包中有 \(A\to\alpha\cdot B\beta\)，且文法中有 \(B\to\gamma\)，则加入 \(B\to\cdot\gamma\)。
  \item 重复上述过程，直到没有新 item 可加入。
\end{enumerate}
\end{defbox}

通俗解释：点号后面如果是非终结符 \(B\)，说明下一步可能要识别一个 \(B\)。既然要识别 \(B\)，就必须把 \(B\) 的所有候选产生式都放进当前状态。

\begin{defbox}[GOTO]
\(\GOTO(I,X)\) 表示从 item 集 \(I\) 通过文法符号 \(X\) 能到达的新 item 集。具体做法是：找到所有 \(A\to\alpha\cdot X\beta\)，把点号移过 \(X\) 得到 \(A\to\alpha X\cdot\beta\)，再对这些 item 取闭包。
\end{defbox}

构造 LR(0) 状态集的算法：

\begin{enumerate}
  \item 构造增广文法 \(G'\)。
  \item 初始状态为 \(\CLOSURE(\{S'\to\cdot S\})\)。
  \item 对每个已有状态 \(I\) 和每个文法符号 \(X\)，计算 \(\GOTO(I,X)\)。
  \item 若结果非空且是新状态，加入状态集。
  \item 重复直到不再产生新状态。
\end{enumerate}

所有状态构成规范 LR(0) 项集族（canonical LR(0) collection），它本质上就是识别活前缀的 DFA。

\subsection{例子：\texorpdfstring{\(S\to BB,\ B\to aB\mid b\)}{S -> BB, B -> aB | b} \pg{69--78}}

增广文法：
\[
(0)\ S'\to S,\qquad (1)\ S\to BB,\qquad (2)\ B\to aB,\qquad (3)\ B\to b.
\]
初始状态：
\[
I_0=\{S'\to\cdot S,\ S\to\cdot BB,\ B\to\cdot aB,\ B\to\cdot b\}.
\]

若从 \(I_0\) 经过 \(B\)：
\[
\GOTO(I_0,B)=\CLOSURE(\{S\to B\cdot B\})
=\{S\to B\cdot B,\ B\to\cdot aB,\ B\to\cdot b\}.
\]
若从 \(I_0\) 经过 \(a\)：
\[
\GOTO(I_0,a)=\CLOSURE(\{B\to a\cdot B\})
=\{B\to a\cdot B,\ B\to\cdot aB,\ B\to\cdot b\}.
\]
若从 \(I_0\) 经过 \(b\)：
\[
\GOTO(I_0,b)=\{B\to b\cdot\}.
\]

课件给出的 LR(0) 表可概括为：

\begin{center}
\small
\begin{tabular}{c|ccc|cc}
\toprule
State & \(a\) & \(b\) & \(\$\) & \(S\) & \(B\) \\
\midrule
0 & \(s3\) & \(s4\) &  & 1 & 2 \\
1 &  &  & acc &  &  \\
2 & \(s3\) & \(s4\) &  &  & 5 \\
3 & \(s3\) & \(s4\) &  &  & 6 \\
4 & \(r3\) & \(r3\) & \(r3\) &  &  \\
5 & \(r1\) & \(r1\) & \(r1\) &  &  \\
6 & \(r2\) & \(r2\) & \(r2\) &  &  \\
\bottomrule
\end{tabular}
\end{center}

\subsection{由 DFA 构造 LR(0) 表 \pg{77--80}}

构造规则：

\begin{itemize}
  \item 若状态 \(S_i\) 中有 \(A\to\alpha\cdot a\beta\)，其中 \(a\) 是终结符，并且 \(\GOTO(S_i,a)=S_j\)，则 \(\ACTION[i,a]=s_j\)。
  \item 若状态 \(S_i\) 中有完成项 \(A\to\alpha\cdot\)，且该产生式是第 \(j\) 条，则对所有终结符 \(a\) 填 \(\ACTION[i,a]=r_j\)。
  \item 若状态中有 \(S'\to S\cdot\)，则 \(\ACTION[i,\$]=acc\)。
  \item 若 \(\GOTO(S_i,A)=S_j\)，其中 \(A\) 是非终结符，则 \(\GOTO[i,A]=j\)。
\end{itemize}

LR(0) 的“0”表示分析器不能看未来输入，只能根据当前状态决定动作。因此状态必须非常干净：要么只移入，要么只规约，不能混杂。

\subsection{LR(0) 冲突与局限 \pg{81--95}}

LR(0) 冲突条件：

\begin{itemize}
  \item Reduce-reduce conflict：某状态同时有两个完成项，例如 \(X\to\beta\cdot\) 和 \(Y\to\omega\cdot\)。
  \item Shift-reduce conflict：某状态同时有完成项和可移入项，例如 \(X\to\beta\cdot\) 和 \(Y\to\omega\cdot t\)。
\end{itemize}

课件示例中，文法：
\[
E'\to E,\quad
E\to E+T\mid T\mid V=E,\quad
T\to (E)\mid i\mid i[E],\quad
V\to i
\]
在读入 \(i\) 后可能出现：
\[
T\to i\cdot,\qquad T\to i\cdot [E],\qquad V\to i\cdot.
\]
这里既有完成项，又有点后为 \([ \) 的可移入项，并且有多个完成项，LR(0) 无法决定。

\begin{warnbox}[LR(0) 的根本局限]
LR(0) 看到完成项就倾向于规约，但当前输入符号可能说明“不该现在规约”。因此需要引入 lookahead。SLR(1) 用 FOLLOW 集做粗粒度限制；LR(1) 把精确 lookahead 放进 item。
\end{warnbox}

\section{SLR(1)：用 FOLLOW 限制规约 \pg{96--106}}

\subsection{SLR(1) 的基本思想 \pg{98--99}}

\begin{defbox}[SLR(1)]
SLR(1) 即 Simple LR with 1-token lookahead。它使用与 LR(0) 相同的配置集、表结构和分析算法，但在填 reduce 动作时加入 FOLLOW 集限制：只有当当前输入符号属于 \(\FOLLOW(A)\) 时，才允许用 \(A\to\alpha\) 规约。
\end{defbox}

如果状态中有完成项 \(A\to\alpha\cdot\)，LR(0) 会把规约动作填到所有终结符列；SLR(1) 只把规约动作填到 \(\FOLLOW(A)\) 对应的列。这样很多无意义规约会被删除。

课件示例中：
\[
\FOLLOW(T)=\{+,),],\$\},\qquad \FOLLOW(V)=\{=\}.
\]
当读入 \(id\) 后：

\begin{itemize}
  \item 若下一个输入符号在 \(\FOLLOW(T)\) 中，规约为 \(T\)。
  \item 若下一个输入符号在 \(\FOLLOW(V)\) 中，规约为 \(V\)。
  \item 若下一个输入符号是 \([ \)，则移入，继续识别 \(id[E]\)。
\end{itemize}

\subsection{SLR(1) 文法判定条件 \pg{100}}

一个文法是 SLR(1)，要求每个配置集满足：

\begin{enumerate}
  \item 若集合中有 \(A\to u\cdot xv\)，其中 \(x\) 是终结符，则集合中不能存在完成项 \(B\to w\cdot\) 且 \(x\in\FOLLOW(B)\)。这对应无 shift-reduce conflict。
  \item 若集合中有两个完成项 \(A\to u\cdot\) 和 \(B\to v\cdot\)，则 \(\FOLLOW(A)\cap\FOLLOW(B)=\emptyset\)。这对应无 reduce-reduce conflict。
\end{enumerate}

通俗地说，如果一个状态中有多个可能动作，只要一个输入符号能唯一告诉我们该选哪一个，SLR(1) 就能处理。

\subsection{SLR(1) 的优势与局限 \pg{101--106}}

优势：

\begin{itemize}
  \item 只加一个 lookahead 和 FOLLOW 集，就大幅扩展了 LR(0) 能处理的文法范围。
  \item 仍然使用 LR(0) 项集，所以状态数量不增加。
  \item 同一状态可以同时有移入项和规约项，也可以有多个完成项，只要 FOLLOW 能把动作区分开。
\end{itemize}

局限：

\begin{itemize}
  \item FOLLOW 是全局信息，只说明“某个非终结符在任何合法句子中可能跟随什么符号”。
  \item 当前状态只描述栈顶这段串，不能充分记录栈下方的上下文。
  \item 因此有时 \(\FOLLOW(A)\) 包含某符号，但在当前这个具体上下文中并不应该用 \(A\) 规约。
\end{itemize}

典型限制例子：
\[
S\to L=R\mid R,\qquad L\to *R\mid id,\qquad R\to L.
\]
输入 \(id=id\) 时，读入第一个 \(id\) 后可规约为 \(L\)。在某状态中同时存在：
\[
S\to L\cdot =R,\qquad R\to L\cdot.
\]
看到下一个输入符号 \(=\) 时，正确动作应是 shift，因为要继续识别赋值式 \(L=R\)。但由于 \(=\in\FOLLOW(R)\)，SLR(1) 也会允许 \(R\to L\) 规约，于是出现 shift-reduce conflict。

这说明：只靠 FOLLOW 不够，需要更精确地记录“当前完成项是在什么上下文里出现的”，这就是 LR(1) 的动机。

\section{LR(1)：把 Lookahead 放进 Item \pg{107--122}}

\subsection{LR(1) item \pg{107--108}}

\begin{defbox}[LR(1) item]
LR(1) item 的一般形式为：
\[
[A\to X_1\cdots X_i\cdot X_{i+1}\cdots X_j,\ a].
\]
其中点号左边表示已经识别的部分，点号右边表示期待识别的部分，\(a\) 是 lookahead。若 item 已完成：
\[
[A\to X_1\cdots X_j\cdot,\ a],
\]
则只有当下一个输入符号正好是 \(a\) 时，才允许规约。
\end{defbox}

对 reduce 条件的精度可以这样比较：

\begin{center}
\small
\begin{tabular}{lll}
\toprule
分析器 & 规约条件 & 精度 \\
\midrule
LR(0) & 状态中有完成项就规约 & 最粗 \\
SLR(1) & 下一输入符号在 \(\FOLLOW(A)\) 中才用 \(A\to\alpha\) 规约 & 中等 \\
LR(1) & 下一输入符号等于 item 的 lookahead 才规约 & 最精确 \\
\bottomrule
\end{tabular}
\end{center}

LR(1) 的代价是状态变多：很多 LR(0) item 的 core 相同，但 lookahead 不同，因此会被区分为不同 LR(1) item，进而产生更多状态。

\subsection{LR(1) 的 CLOSURE 和 GOTO \pg{109--110}}

LR(1) 状态构造与 LR(0) 类似，但闭包和跳转要传播 lookahead。

\begin{defbox}[LR(1) CLOSURE]
若 \(I\) 中有 item \([A\to u\cdot Bv,\ a]\)，且文法有 \(B\to w\)，则对每个
\[
b\in\FIRST(va)
\]
加入 item \([B\to\cdot w,\ b]\)。
\end{defbox}

为什么是 \(\FIRST(va)\)？因为当前正在期望一个 \(B\)，而 \(B\) 后面紧跟 \(v\)，整个 \(A\) 规约后又可能遇到 lookahead \(a\)。因此 \(B\) 规约完成后可能看到的符号就是 \(va\) 能推出串的首符号。

\begin{defbox}[LR(1) GOTO]
若 \(I\) 中有 \([A\to u\cdot Xv,\ a]\)，则
\[
\GOTO(I,X)=\CLOSURE(\{[A\to uX\cdot v,\ a]\}).
\]
和 LR(0) 一样移动点号，但必须携带并传播 lookahead。
\end{defbox}

构造步骤：

\begin{enumerate}
  \item 从 \(\CLOSURE([S'\to\cdot S,\ \$])\) 开始。
  \item 对每个状态和文法符号计算 \(\GOTO\)。
  \item 重复直到没有新状态。
\end{enumerate}

\subsection{LR(1) 示例 \pg{109--121}}

文法：
\[
(0)\ S'\to S,\quad (1)\ S\to XX,\quad (2)\ X\to aX,\quad (3)\ X\to b.
\]
初始状态的一部分：
\[
\begin{aligned}
I_0=\{&S'\to\cdot S,\ \$;\\
&S\to\cdot XX,\ \$;\\
&X\to\cdot aX,\ a/b;\\
&X\to\cdot b,\ a/b\}.
\end{aligned}
\]
这里 \(X\) 的 lookahead 是 \(a/b\)，因为在 \(S\to\cdot XX\) 中，第一个 \(X\) 后面还有一个 \(X\)，而 \(\FIRST(X\$)=\{a,b\}\)。

当进入 \(S\to X\cdot X,\ \$\) 后，第二个 \(X\) 后面直接是 \(\$\)，所以新加入的 \(X\) 相关 item 的 lookahead 是 \(\$\)：
\[
X\to\cdot aX,\ \$,\qquad X\to\cdot b,\ \$.
\]

这正是 LR(1) 比 SLR(1) 更精确的地方：同样是 \(X\to b\cdot\)，在不同状态中可能只在 \(a/b\) 下规约，或只在 \(\$\) 下规约。

\subsection{LR(1) 解析表与文法条件 \pg{120--122}}

LR(1) 表构造：

\begin{itemize}
  \item 若 \([A\to\alpha\cdot a\beta,\ b]\in S_i\) 且 \(\GOTO(S_i,a)=S_j\)，则 \(\ACTION[i,a]=s_j\)。Shift 不关心 lookahead。
  \item 若 \([A\to\alpha\cdot,\ a]\in S_i\)，则只在 \(\ACTION[i,a]\) 填入 reduce。
  \item \(\GOTO\) 表仍由非终结符转移决定。
\end{itemize}

一个文法是 LR(1)，要求每个 LR(1) 配置集满足：

\begin{enumerate}
  \item 对任意移入项 \([A\to u\cdot xv,\ a]\)，不存在完成项 \([B\to v'\cdot,\ x]\)。这表示无 shift-reduce conflict。
  \item 所有完成项的 lookahead 两两不冲突，即不能同时有 \([A\to u\cdot,\ a]\) 和 \([B\to v\cdot,\ a]\)。这表示无 reduce-reduce conflict。
\end{enumerate}

每个 SLR(1) 文法都是 LR(1)，但 LR(1) parser 可能有更多状态，因为 LR(1) 会按不同 lookahead 拆分状态。

\section{LALR(1)：合并同 Core 状态的折中方案 \pg{123--139}}

\subsection{为什么需要 LALR(1) \pg{124}}

LR(1) 的缺点是状态数可能远多于 SLR(1) 或 LR(0)。一个 LR(0) item 可能因为不同 lookahead 被拆成多个 LR(1) item，理论上可能带来指数级状态增长。

\begin{defbox}[LALR(1)]
LALR(1) 即 lookahead LR。它先获得 LR(1) 的精确 lookahead 信息，再合并“core 相同”的状态，以减少状态数。它试图在 LR(1) 的表达能力和 SLR(1) 的表规模之间折中。
\end{defbox}

\begin{defbox}[Core]
LR(1) item 去掉 lookahead 后剩下的 LR(0) item 称为 core。若两个状态中的 item 数量相同，且所有 item 的 core 相同，只是 lookahead 不同，则这些状态可以尝试合并。
\end{defbox}

\subsection{状态合并示例 \pg{125--127}}

课件中，LR(1) 状态：
\[
\begin{aligned}
I_3 &: X\to a\cdot X,\ a/b;\ X\to\cdot aX,\ a/b;\ X\to\cdot b,\ a/b,\\
I_6 &: X\to a\cdot X,\ \$;\ X\to\cdot aX,\ \$;\ X\to\cdot b,\ \$.
\end{aligned}
\]
二者 core 相同，只是 lookahead 不同，因此合并为：
\[
I_{36}: X\to a\cdot X,\ a/b/\$;\quad X\to\cdot aX,\ a/b/\$;\quad X\to\cdot b,\ a/b/\$.
\]
类似地，\(I_4\) 与 \(I_7\) 合并为 \(I_{47}\)，\(I_8\) 与 \(I_9\) 合并为 \(I_{89}\)。合并后表的行数减少。

\subsection{合并效果：冲突与错误延迟 \pg{128--133}}

合并可能产生两个效果：

\begin{enumerate}
  \item 可能引入 reduce-reduce conflict。
  \item 不会新引入 shift-reduce conflict。
  \item 错误检测可能延迟：LALR 不会比 LR 多移入错误符号，但可能先做更多规约才发现错误。
\end{enumerate}

为什么不会新引入 shift-reduce conflict？如果合并后的状态中出现完成项 \([A\to\alpha\cdot,a]\) 和移入项 \([B\to\beta\cdot a\gamma,b]\)，由于合并要求 core 相同，原始某个 LR(1) 状态中也会包含对应 core 的移入项，因此 shift-reduce conflict 本来就存在，不能说是合并新造成的。

但 reduce-reduce conflict 可能新出现。例如：
\[
S\to aBc\mid bCc\mid aCd\mid bBd,\qquad B\to e,\qquad C\to e.
\]
两个 LR(1) 状态原本可能分别在不同 lookahead 下规约 \(B\to e\) 或 \(C\to e\)，合并后 lookahead 集合并，可能在同一输入符号下同时可规约为 \(B\) 和 \(C\)。

\subsection{LALR 表构造与文法关系 \pg{134--137}}

构造 LALR 表有两种思路：

\begin{itemize}
  \item 直接方法：先构造完整 LR(1) 状态，再合并同 core 状态。简单但开销大。
  \item 高效方法：构造状态时就即时合并，避免完整展开巨大的 LR(1) 自动机。
\end{itemize}

LALR(1) 表的状态数通常与 LR(0)、SLR(1) 相同，但规约动作比 SLR(1) 更少，因为 lookahead 更精确。

文法包含关系：
\[
LR(1) \supset LALR(1) \supset SLR(1) \supset LR(0).
\]
课件总结：

\begin{itemize}
  \item 支持文法范围：\(LR > LALR > SLR > LR(0)\)。
  \item 状态数量：\(LR > LALR = SLR = LR(0)\)。
  \item 若文法是 SLR(1)，它一定也是 LALR(1) 和 LR(1)。
  \item 若文法是 LR(1)，不一定是 LALR(1)，因为合并可能引入 reduce-reduce conflict。
  \item 若文法不是 SLR(1)，仍可能是 LALR(1)，因为更精确的 lookahead 可能消除 SLR 冲突。
\end{itemize}

\subsection{LL vs LR \pg{138--139}}

LL(k) 和 LR(k) 的核心差别在于做选择的时机。

\begin{longtable}{L{0.19\linewidth}L{0.36\linewidth}L{0.35\linewidth}}
\toprule
对比项 & LL(k) & LR(k) \\
\midrule
分析方向 & 从根到叶，预测推导 & 从叶到根，反向规约 \\
选择依据 & 栈顶非终结符 \(A\) 与 RHS 开头 \(k\) 个终结符 & 栈顶已完整识别的 RHS 与其后 \(k\) 个 lookahead \\
选择时机 & 还没看到完整 RHS 就要猜产生式 & 等整个 RHS 已经在栈顶后再决定规约 \\
左因子 & 公共前缀会造成预测困难 & 可以等公共前缀之后的信息出现 \\
左递归 & 递归下降不适合直接处理 & 可以自然处理 \\
能力 & 较弱 & 较强，课件写作 \(LL<LR\) \\
\bottomrule
\end{longtable}

通俗地说，LL 像“开头看几眼就猜路”，LR 像“先把一整段证据收集完，再决定它是哪条产生式的右部”。所以 LR 能处理更多文法。

\section{Syntax Analysis Revisit：语法分析总复习 \pg{140--189}}

课件后半段以 Revisit 形式复习前几讲。下面按其顺序补齐知识点。

\subsection{语法分析的位置与文法定义 \pg{141--145}}

语法分析是编译第二阶段：

\begin{itemize}
  \item 输入：词法分析产生的 token 序列。
  \item 输出：分析树（parse tree）或抽象语法树（AST）。
\end{itemize}

正则表达式和有限自动机不足以表达嵌套结构，例如 \(\{a^n b^n\mid n>0\}\)，因此需要文法，尤其是上下文无关文法 CFG。

文法四元组：
\[
G=(V_T,V_N,S,\delta).
\]

\begin{longtable}{L{0.18\linewidth}L{0.72\linewidth}}
\toprule
组成 & 含义 \\
\midrule
\(V_T\) & 终结符集合，来自词法分析的 token，通常是 parse tree 的叶子。\\
\(V_N\) & 非终结符集合，表示语法结构，通常是 parse tree 的内部结点。\\
\(S\) & 开始符号。\\
\(\delta\) & 产生式集合，形式为 LHS \(\to\) RHS。\\
\bottomrule
\end{longtable}

推导是从 LHS 到 RHS 应用产生式，从开始符号生成句子；规约是推导的逆过程，从输入串回到开始符号。分析树表示推导或规约的层次结构，但隐藏了具体应用产生式的先后顺序。

\subsection{句型、句子、语言、二义性 \pg{144--145}}

以文法 \(A\to c\mid Ab\) 为例：
\[
A\Rightarrow c,\quad
A\Rightarrow Ab\Rightarrow cb,\quad
A\Rightarrow Abb\Rightarrow cbb,\ldots
\]
可生成语言 \(\{cb^n\mid n\ge 0\}\)。

\begin{itemize}
  \item 句型：从开始符号经过若干步推导得到的符号串，可含非终结符。
  \item 句子：只含终结符的句型。
  \item 语言：文法可生成的所有句子的集合。
  \item 二义文法：某个句子有不止一棵分析树，等价地有不止一个最左或最右推导。
\end{itemize}

二义性可通过定义优先级、结合性或改写文法来消除。

\subsection{乔姆斯基文法体系与 AST \pg{146--148}}

乔姆斯基将文法分为四类：

\begin{center}
\small
\begin{tabular}{lll}
\toprule
类型 & 名称 & 典型能力 \\
\midrule
Type 0 & 无限制文法 & 图灵机可识别语言 \\
Type 1 & 上下文有关文法 & 上下文有关语言 \\
Type 2 & 上下文无关文法 CFG & 程序语言语法主体 \\
Type 3 & 正则文法 & 正则语言，可由 FA 识别 \\
\bottomrule
\end{tabular}
\end{center}

AST 是 parse tree 的简化表示。它删除不影响语义的推导细节，例如只用于表达优先级和括号结构的中间非终结符，让后续语义分析和代码生成更方便。

\subsection{Top-down 复习：RDP、FIRST、FOLLOW 与 LL(1) \pg{149--160}}

Top-down parser 从根到叶构造语法树，核心问题是：

\begin{itemize}
  \item 当前替换哪个非终结符？
  \item 选择该非终结符的哪条产生式？
\end{itemize}

递归下降分析（RDP）直观，但存在两个问题：

\begin{itemize}
  \item 左递归会导致无限递归，需要消除。
  \item 回溯开销大且实现复杂，应尽量避免。
\end{itemize}

直接左递归消除公式：
\[
A\to A\alpha_1\mid\cdots\mid A\alpha_m\mid \beta_1\mid\cdots\mid\beta_n
\]
改写为：
\[
A\to \beta_1A'\mid\cdots\mid\beta_nA',\qquad
A'\to \alpha_1A'\mid\cdots\mid\alpha_mA'\mid\eps.
\]

FIRST 与 FOLLOW：

\begin{itemize}
  \item \(\FIRST(\alpha)\)：\(\alpha\) 能推出的句子的首个终结符集合，可包含 \(\eps\)。
  \item \(\FOLLOW(A)\)：在某个句型中能紧跟在 \(A\) 后面的终结符集合，可包含结束符 \(\$\)。
\end{itemize}

表达式文法消左递归后：
\[
\begin{aligned}
E&\to TE', & E'&\to +TE'\mid\eps,\\
T&\to FT', & T'&\to *FT'\mid\eps,\\
F&\to (E)\mid id.
\end{aligned}
\]
课件给出的集合为：
\[
\begin{aligned}
\FIRST(E)&=\{(,id\}, & \FIRST(T)&=\{(,id\},\\
\FIRST(E')&=\{+,\eps\}, & \FIRST(T')&=\{*,\eps\},\\
\FIRST(F)&=\{(,id\},\\
\FOLLOW(E)&=\{\$,)\}, & \FOLLOW(T)&=\{+,\$,)\},\\
\FOLLOW(E')&=\{\$,)\}, & \FOLLOW(T')&=\{+,\$,)\},\\
\FOLLOW(F)&=\{*,+,\$,)\}.
\end{aligned}
\]

LL(1) 分析表的行是非终结符，列是终结符加 \(\$\)。非递归 LL(1) 算法根据栈顶符号 \(X\) 和当前输入符号 \(a\) 决定动作：

\begin{itemize}
  \item 若 \(X\) 是终结符且 \(X=a\)，匹配并前进。
  \item 若 \(X\) 是非终结符且 \(M[X,a]\) 有产生式，弹出 \(X\)，把 RHS 逆序压栈。
  \item 若不匹配或表项为空，报错。
  \item 若 \(X=a=\$\)，接受。
\end{itemize}

\subsection{Bottom-up 复习 \pg{161--189}}

Revisit 最后再次串联本讲主线：

\begin{itemize}
  \item Shift-reduce parser 从输入规约到开始符号，是最右推导逆过程。
  \item Phrase、simple phrase、handle 用于说明“该规约哪里”。
  \item Viable prefix 说明 LR 栈内容始终没有越过句柄右端。
  \item LR(k) parser 用状态序列和 \ACTION/\GOTO 表驱动分析。
  \item LR(0) 用 item、\CLOSURE、\GOTO 构造自动机，但不看 lookahead，能力最弱。
  \item SLR(1) 用 FOLLOW 限制 reduce。
  \item LR(1) 在 item 中加入 lookahead，精确控制 reduce。
  \item LALR(1) 合并 LR(1) 同 core 状态，常用于实践。
\end{itemize}

\section{练习题与解答 \pg{190--206}}

\subsection*{题 1：正则表达式能否表达有限语言和嵌套计数语言？}

判断正则表达式是否可以表达：
\[
L_1=\{x^n y z^n\mid 1\le n\le 10\},\qquad
L_2=\{x^n y z^n\mid n>0\}.
\]

\textbf{解答：} \(L_1\) 可以。因为 \(n\) 只从 1 到 10，语言是有限集；任何有限语言都是正则语言，都可写成有限个串的并。\(L_2\) 不可以。它要求 \(x\) 与 \(z\) 的数量任意大且相等，有限自动机没有无限计数能力，因此不是正则语言。

\subsection*{题 2：DFA 没有 epsilon 弧，能识别空串吗？}

\textbf{解答：} 可以。DFA 没有 \(\eps\) 转移不代表不能接受 \(\eps\)。当且仅当初始状态本身是接受状态时，DFA 接受空串。

\subsection*{题 3：Thompson 构造 \texorpdfstring{\((a\mid b)^*\)}{(a|b)*} 有几个状态？}

\textbf{解答：} 8 个。按 Thompson 构造，\(a\) 与 \(b\) 各 2 个状态，选择结构增加入口/出口，Kleene star 再增加入口/出口，课件答案为 8。

\subsection*{题 4：无二义文法的最左推导和最右推导分析树是否一样？}

\textbf{解答：} 一样。无二义文法中，一个句子只有一棵分析树。最左推导和最右推导只是应用产生式的顺序不同，分析树会隐去推导顺序，因此得到同一棵树。

\subsection*{题 5：如何判断文法是 LL(1)？}

对每个非终结符的候选式 \(A\to\alpha\mid\beta\)，需要满足：
\[
\FIRST(\alpha)\cap\FIRST(\beta)=\emptyset.
\]
若某个候选式可推出 \(\eps\)，例如 \(\beta\Rightarrow^*\eps\)，还要满足：
\[
\FIRST(\alpha)\cap\FOLLOW(A)=\emptyset.
\]

通俗解释：看一个输入符号时，不能同时指向两条候选式；如果某条候选式能为空，则当输入符号属于 FOLLOW 时才选空产生式，也不能和其他候选式冲突。

\subsection*{题 6：LL(1) 表的行数和列数由什么决定？}

\textbf{解答：} 行数由非终结符个数决定；列数由终结符个数加 1 决定，额外的 1 是结束符号 \(\$\)。

\subsection*{题 7：活前缀题}

设句型 \(aaAb\) 的最右推导为：
\[
S\Rightarrow aBb\Rightarrow aaAb,
\]
句柄为 \(aA\)。该句型的活前缀是什么？

\textbf{解答：} 活前缀不能越过句柄 \(aA\) 的右端，因此为：
\[
a,\quad aa,\quad aaA.
\]

\subsection*{题 8：找短语、直接短语与句柄}

文法：
\[
E\to T\mid E+T,\quad T\to F\mid T*F,\quad F\to (E)\mid a\mid b\mid c.
\]
句型：\(a*b+c\)。

\textbf{解答：}
\[
\text{短语： } a*b+c,\ a*b,\ a,\ b,\ c.
\]
\[
\text{直接短语： } a,\ b,\ c.
\]
\[
\text{句柄： } a.
\]
理由是这些短语必须对应语法树中的完整子树叶子序列，而句柄是最左直接短语。

\subsection*{题 9：构造 LL(1) 分析表}

文法：
\[
S\to A,\qquad A\to\eps,\qquad A\to bbA.
\]

\textbf{解答：}
\[
\FIRST(A)=\{b,\eps\},\qquad \FOLLOW(S)=\{\$\},\qquad \FOLLOW(A)=\{\$\}.
\]
分析表：

\begin{center}
\begin{tabular}{c|cc}
\toprule
 & \(b\) & \(\$\) \\
\midrule
\(S\) & \(S\to A\) & \(S\to A\) \\
\(A\) & \(A\to bbA\) & \(A\to\eps\) \\
\bottomrule
\end{tabular}
\end{center}

解释：\(S\to A\) 的 FIRST 含 \(b\) 和 \(\eps\)，因此在 \(b\) 和 FOLLOW(S) 的 \(\$\) 列填表；\(A\to bbA\) 在 \(b\) 列，\(A\to\eps\) 在 FOLLOW(A) 的 \(\$\) 列。

\subsection*{题 10：给定文法构造 LR(1) DFA 和分析表}

文法：
\[
(1)\ S\to A,\qquad
(2)\ A\to A+A,\qquad
(3)\ A\to B++,\qquad
(4)\ B\to y.
\]
课件给出的 LR(1) 项集核心如下：

\[
\begin{aligned}
I_0:&\ S\to\cdot A,\ \$;\ A\to\cdot A+A,\ +/\$;\ A\to\cdot B++,\ +/\$;\ B\to\cdot y,\ +.\\
I_1:&\ B\to y\cdot,\ +.\\
I_2:&\ S\to A\cdot,\ \$;\ A\to A\cdot +A,\ +/\$.\\
I_3:&\ A\to B\cdot ++,\ +/\$.\\
I_4:&\ A\to A+\cdot A,\ +/\$;\ A\to\cdot A+A,\ +/\$;\ A\to\cdot B++,\ +/\$;\ B\to\cdot y,\ +.\\
I_5:&\ A\to A+A\cdot,\ +/\$;\ A\to A\cdot +A,\ +/\$.\\
I_6:&\ A\to B+\cdot +,\ +/\$.\\
I_7:&\ A\to B++\cdot,\ +/\$.
\end{aligned}
\]

对应表中关键动作：

\begin{center}
\small
\begin{tabular}{c|ccc|cc}
\toprule
State & \(y\) & \(+\) & \(\$\) & \(A\) & \(B\) \\
\midrule
0 & \(s1\) &  &  & 2 & 3 \\
1 &  & \(r4\) &  &  &  \\
2 &  & \(s4\) & acc &  &  \\
3 &  & \(s6\) &  &  &  \\
4 & \(s1\) &  &  & 5 & 3 \\
5 &  & \(r2/s4\) & \(r2\) &  &  \\
6 &  & \(s7\) &  &  &  \\
7 &  & \(r3\) & \(r3\) &  &  \\
\bottomrule
\end{tabular}
\end{center}

\subsection*{题 11：该文法是否为 LR(1)？如何解决？}

\textbf{解答：} 不是 LR(1)。从表中可见状态 5 在输入 \(+\) 时存在 \(r2/s4\) 冲突：既可以按 \(A\to A+A\) 规约，也可以移入 \(+\) 继续形成更长的 \(A+A+A\)。

可选解决方法：

\begin{itemize}
  \item 改写文法，明确 \(+\) 的结合性和优先级。
  \item 在分析器生成工具中声明冲突解决策略，例如规定 \(+\) 左结合时倾向规约，右结合时倾向移入。
  \item 如果语言设计允许，也可修改语法结构避免这种二义选择。
\end{itemize}

\subsection*{题 12：分析句子 \texorpdfstring{\(y++\)}{y++}}

使用上一题文法与表，分析 \(y++\)：

\begin{center}
\small
\begin{tabular}{llll}
\toprule
状态栈 & 符号栈 & 输入 & 动作/跳转 \\
\midrule
0 & \(\$\) & \(y++\$\) & \(s1\) \\
0 1 & \(\$y\) & \(++\$\) & \(r4: B\to y\) \\
0 & \(\$B\) & \(++\$\) & \(\GOTO(0,B)=3\) \\
0 3 & \(\$B\) & \(++\$\) & \(s6\) \\
0 3 6 & \(\$B+\) & \(+\$\) & \(s7\) \\
0 3 6 7 & \(\$B++\) & \(\$\) & \(r3: A\to B++\) \\
0 & \(\$A\) & \(\$\) & \(\GOTO(0,A)=2\) \\
0 2 & \(\$A\) & \(\$\) & acc \\
\bottomrule
\end{tabular}
\end{center}

因此 \(y++\) 是该文法的一个句子。

\subsection*{题 13：LR(0)、SLR(1)、LR(1)、LALR(1) 的关系}

\textbf{解答：}
\[
LR(1) \supset LALR(1) \supset SLR(1) \supset LR(0).
\]
能力越往左越强。状态数一般为：
\[
LR(1) > LALR(1)=SLR(1)=LR(0).
\]
LR(0) 不看展望符号；SLR(1) 用 FOLLOW 限制规约；LR(1) 用 item 自带 lookahead 精确规约；LALR(1) 合并 LR(1) 同 core 状态，在能力和表规模之间折中。

\section*{复习建议}
\addcontentsline{toc}{section}{复习建议}

本讲复习时建议按三层掌握：

\begin{enumerate}
  \item 概念层：能解释 handle、viable prefix、item、closure、goto、lookahead、core。
  \item 机制层：能手工走一遍 shift-reduce 分析，能根据 \ACTION/\GOTO 表分析输入串。
  \item 构造层：能从小文法构造 LR(0) 项集和表，能说明 SLR(1)、LR(1)、LALR(1) 分别如何改进 LR(0)。
\end{enumerate}

最容易混淆的点是：FOLLOW 是 SLR(1) 的粗粒度全局信息；LR(1) item 的 lookahead 是带上下文的局部信息；LALR(1) 则把 LR(1) 的局部信息合并到同 core 状态中，所以表小但可能产生 reduce-reduce conflict。

\end{document}
